\section{Original-level Type II exclusion on the quadratic branch}
\label{sec:cubic-original-typeii-exclusion}

The original \(N_\psi=4\), \(\psi=\chi_{-4}\) Type II point set at \(269\) is empty
under the exact old modularity and local inputs retained in
AC-G68-point.  The proof treats every weight allowed by the literal
original datum.  It does not assert a nonempty seed satisfying those
same original local conditions, nor does it close H04.4.

Put
\[
 p=269,\qquad K=\mathbb Q(\sqrt{-3}),\qquad
 \eta=\eta_K=\chi_{-3},\qquad \psi=\chi_{-4},\qquad
 e\ge1,\qquad268\nmid2e,\qquad k=2e+1.
 \tag{G116.1}
\]
The symbol \(c_p\) means the actual action on \(p\)-power roots
of unity.  At \(q\ne p\), arithmetic Frobenius \(\sigma_q\) has
\(c_p(\sigma_q)=q\); geometric Frobenius is
\(\operatorname{Fr}_q=\sigma_q^{-1}\).
The source \texttt{Newform.Att} evaluates its representation at
\(\sigma_q^{-1}\).  Consequently its modular representation for
the old normalized endpoint \(V\) is
\[
                   R_f=V c_p^{-e}.
 \tag{G116.2}
\]
Its geometric Frobenius polynomial at a good prime is
\(X^2-a_q(f)X+\eta(q)q^{2e}\).
This fixes the normalization throughout the argument.

\begin{theorem}[AC-G116-exclusion: no original quadratic-character Type II points]
\label{thm:ACG116exclusion}
Retain the literal original datum with conductor modulus
\(N_\psi=4\), the actual cubic ring and selected local prime of
AC-G65, and every old automorphy and local hypothesis used in
AC-G68-point, including the original coefficient condition
\(e_O=1\), \texttt{GD}, \texttt{CohInp}, and
\texttt{AutomorphyUnramifiedCoefficientsAlt}.
For every weight in (G116.1), there is no actual characteristic-zero
Type II point of that cubic ring with the original
\(N_\psi=4\) local conditions.

More precisely, the contradiction uses the following consequences
of such a point, all supplied by AC-G68-block and AC-G68-point:
an integral irreducible odd \(V:G_{\mathbb Q}\to\operatorname{GL}_2(E)\),
for a finite coefficient field \(E\), with determinant \(\eta\),
residual semisimplification
\[
 \overline V^{\mathrm{ss}}=
       (\eta\psi\overline c_p^{\,e})
          \oplus(\psi^{-1}\overline c_p^{-e}),
 \tag{G116.3}
\]
the exact original away-inertia condition, the ordered ordinary
potentially crystalline condition at \(p\), and an actual cuspidal
newform attached to (G116.2).  The published classical local--global
compatibility and newform theorems used below are retained inputs.
The proof excludes any \(V\) with these displayed residual,
ordinary, exact-inertia and modular properties, regardless of how
it was obtained.  In particular it applies directly to a native
old Type II endpoint with those properties; no embedding of an
arbitrary old point into the selected cubic component is assumed.
\end{theorem}

\begin{proof}
\proofstep{1. The exact original inertia and the primitive twist.}
The literal source \texttt{GD.gen} is \(268\nmid2e\), and
\texttt{GD.par} for \(\psi=\chi_{-4}\) makes \(e\) odd.
The original source is not being replaced by a weakened datum.
For the argument below, the relevant residue of \(k=2e+1\) is
an odd integer \(n\) in \(\{3,5,\ldots,267\}\): residue one
would be equivalent to \(268\mid2e\).  In particular both
\(e\equiv133\) and \(e\equiv267\pmod{268}\) are included.
The remaining local and automorphy conclusions of AC-G68-point
hold for every such retained datum; none assumes a particular
weight residue.

Assume that such a Type II point exists.  Removing its cubic
endpoint character gives the actual old point of AC-G68-point.
In particular, on inertia inside \(G_K\) away from \(p\),
\(V\) is diagonal with characters \(\psi,\psi^{-1}\).
Equivalently, one may use the literal old residual-kernel killing
condition and the prime-to-\(p\) integral character projectors.
The following assertions concern the representation itself,
including zero monodromy, rather than only its residual trace.

At \(2\), \(K/\mathbb Q\) is unramified, so all rational inertia
lies in \(G_K\).  Since \(\psi^{-1}=\psi\), it acts on \(V\)
as the scalar \(\psi\).  At \(3\), inertia over \(K\) acts
trivially.  Rational inertia therefore factors through
\(\operatorname{Gal}(K_3/\mathbb Q_3)=C_2\), and its determinant
is \(\eta\).  Its two characters are exactly \(1,\eta\).
At every other prime away from \(p\), \(V\) is unramified.
These statements also follow directly from the full inertia
argument in \texttt{AdmissibleOfTypeII.lean}; the source's
\texttt{Level} parameter \(N_\psi|D_K|=12\) is not being
identified with the classical conductor of a newform.

At \(p\), \(\psi\) is unramified, and its conductor-primary
component is trivial.  AC-G68-point gives potential crystallinity
and trivial inertial type for \(V\).  Its crystalline descent
argument, over the unramified quadratic \(K_p/\mathbb Q_p\),
therefore gives crystallinity over \(\mathbb Q_p\).
The ordered ordinary line has inertia character \(c_p^e\);
the quotient has inertia character \(c_p^{-e}\).
This is also exactly the good-twist proof of AC-K29-inert,
Step 1, with its conductor-primary character \(\tau=1\).
That proof is used here without any Euler-system assertion.

Let \(\pi_f\) be the cuspidal automorphic representation of the
newform supplied by AC-G68-point.  Define \(g\) to be the
normalized \emph{primitive newform of the automorphic twist}
\[
                       \pi_f\otimes\bigl((\psi\eta)\circ\det\bigr).
 \tag{G116.4}
\]
Twisting preserves cuspidality and weight, and its nebentypus is
still \(\eta\), since \((\psi\eta)^2=1\).
Its attached representation, in convention (G116.2), is
\[
                      R_g=V\psi\eta c_p^{-e}.
 \tag{G116.5}
\]
This identification follows from the good-prime eigenvalues and
Chebotarev, or directly from twisting compatibility.
Both quadratic twists in (G116.4) are primitive automorphic
twists.  Coefficientwise Dirichlet twisting would erase
coefficients at the twisting conductor and does not give the
primitive coefficients used below.

The exact inertia just computed shows that \(R_g\) is unramified
at every prime different from \(3,p\).  At \(3\) its inertia
is \(1\oplus\eta\), with zero monodromy, and hence its Artin
conductor exponent is one.  Indeed \(\eta\) has conductor
\(3\) and is tame there.  At \(p\), \(R_g\) is crystalline.
Deligne--Langlands--Carayol compatibility away from \(p\),
and Saito compatibility at \(p\), identify these full
Frobenius-semisimple Weil--Deligne parameters with those of
the local automorphic factors.  The crystalline parameter at
\(p\) is unramified with zero monodromy, so that factor is
spherical.  Classical newform conductor theory now gives
\[
                    g\in S_k(\Gamma_0(3),\eta)
                 \quad\hbox{of primitive conductor }3.
 \tag{G116.6}
\]
For comparison, the untwisted \(V\) has the scalar character
\(\psi\) of conductor \(4\) twice at \(2\), and hence Artin
conductor exponent four there.  Its primitive modular conductor is \(48\), not
the formal source \texttt{Level} parameter \(12\).
No level-lowering congruence theorem is being invoked: (G116.6)
is the conductor of the actual characteristic-zero twist.

\proofstep{2. Good-prime coefficients and the exact root at \(3\).}
From (G116.3) and (G116.5),
\[
             \overline R_g^{\mathrm{ss}}
                    =1\oplus\eta\overline c_p^{-2e}.
 \tag{G116.7}
\]
Thus for every prime \(q\ne3,p\), geometric Frobenius gives
\[
                    a_q(g)\equiv1+\eta(q)q^{k-1}
                              \pmod{\mathfrak m_E}.
 \tag{G116.8}
\]
The extra \(\eta\) in (G116.4) is essential for this displayed
orientation.  Twisting by \(\psi\) alone would give
\(\eta(q)+q^{2e}\), the opposite Eisenstein ordering.

Choose an inertia element at \(3\) whose image in the quadratic
inertia quotient is nontrivial.  Its action \(s\) on the stable
integral lattice satisfies \(s^2=1\).  Since \(2\) is a unit,
\((1+s)/2\) is an integral rank-one projector onto the actual
inertia-invariant line.  It is stable under the decomposition
group: conjugation acts trivially on the quotient \(C_2\).
Taking this projector commutes with reduction.  In (G116.7),
the unique inertia-trivial constituent at \(3\) is the character
\(1\), so the invariant line's unramified character reduces to
the trivial character on the full decomposition group.

The local Euler-factor identity in the cohomological, geometric
Frobenius convention is
\[
 L_3(g,X)^{-1}
   =\det(1-X\operatorname{Fr}_3\mid R_g^{I_3})
   =1-a_3(g)X.
 \tag{G116.9}
\]
Consequently
\[
                             a_3(g)\equiv1.
 \tag{G116.10}
\]
This uses the primitive newvector Euler factor of the actual
automorphic twist, and not the zero coefficient of a
nonprimitive coefficientwise twist.

\proofstep{3. The ordinary root at the coefficient prime.}
The ordinary line of \(R_g\) is unramified, because its
\(c_p^e\) inertia in \(V\) is cancelled by (G116.5).
Its saturated intersection with the original stable lattice
has rank one.  Its reduction is determined by inertia:
the two characters in (G116.7) restrict to
\(1\) and \(\overline c_p^{-2e}\), respectively, and
\(268\nmid2e\).  Thus the unramified line reduces to the first
character in (G116.7), namely \(1\) on the full decomposition
group.  Write its unramified characteristic-zero character as
\(u\).  Then \(\overline u=1\).

The eigenvalue of crystalline Frobenius on this unramified line
is \(\alpha=u(\operatorname{Fr}_p)\), where any arithmetic
Frobenius lift is inverted.  It is an integral unit and
\(\alpha\equiv1\).  Saito compatibility and the good-prime
Hecke polynomial show that \(\alpha\) is the unit root of
\[
                        X^2-a_p(g)X+\eta(p)p^{2e}.
 \tag{G116.11}
\]
The other root is \(\eta(p)p^{2e}/\alpha\); it is integral and
has positive valuation.  Therefore
\[
                             a_p(g)\equiv1.
 \tag{G116.12}
\]
No residual ordinarity criterion or inversion of \(a_p\)
before identifying this actual line is used.

\proofstep{4. All positive Fourier coefficients.}
Let
\[
 E_k(q)=\frac{L(1-k,\eta)}2+
       \sum_{n\ge1}\left(\sum_{d\mid n}\eta(d)d^{k-1}\right)q^n
 \tag{G116.13}
\]
be the standard Eisenstein series \(E_k(1,\eta)\) of level
\(3\), weight \(k\), and character \(\eta\).
Katz's \(G_{k,\eta}\) has precisely this normalization.
The coefficient at \(q^n\) is denoted \(b_n\).

The primitive newform Euler product gives multiplicativity and,
for \(r\ge1\), the recurrences
\[
 a_{q^{r+1}}=a_q a_{q^r}-\eta(q)q^{k-1}a_{q^{r-1}}
          \quad(q\ne3),\qquad
 a_{3^r}=a_3^r.
 \tag{G116.14}
\]
The numbers \(b_n\) satisfy the same recurrences.
At \(3\), \(b_{3^r}=1\).  At \(p\), \(b_p\equiv1\)
and the second recurrence term in (G116.14) vanishes modulo
\(\mathfrak m_E\).  Equations (G116.8), (G116.10), and
(G116.12) therefore imply, by induction and multiplicativity,
\[
                              a_n(g)\equiv b_n
                                  \quad(n\ge1).
 \tag{G116.15}
\]
All these coefficients are integral.  This follows either from
the integral newform Hecke lattice, or here directly from the
integral Galois traces away from \(p\), the invariant-line
root at \(3\), the two integral roots in (G116.11), and
(G116.14).  Thus no omitted bad Euler factor is concealed in
(G116.15).

\proofstep{5. The integral measure and a finite all-residue certificate.}
For completeness, the required all-weight congruence is derived
from the actual Eisenstein measure rather than an unspecified
Kummer congruence.  Katz, \emph{The Eisenstein measure and p-adic
interpolation}, equations (3.3.1)--(3.3.2) and Theorem 3.3.3,
gives an integral measure \(2H^{a,b}\) on
\(\mathbb Z_p\times\mathbb Z/3\mathbb Z\) with
\[
 H^{a,b}(x^{m-1}\eta(t))
       =\bigl(1-a^m\eta(b)\bigr)G_{m,\eta}
                  \quad(m\ge3\text{ odd}).
 \tag{G116.16}
\]
Take \(a=2\), \(b=2\bmod3\).  Since \(\eta(2)=-1\),
the smoothing factor is \(1+2^m\).
The exact computation \(2^{134}\equiv-1\pmod{269}\)
shows that for every odd \(m\),
\[
             2^m\not\equiv1,-1\pmod{269}.
 \tag{G116.17}
\]
Indeed either equality would have \(134\)th power one, whereas
\((2^m)^{134}=(-1)^m=-1\).  Thus both \(2\) and the smoothing
factor are units, and (G116.16) makes \(G_{m,\eta}\) integral
in the generalized modular-form ring.

Let \(n\in\{3,5,\ldots,267\}\) be the residue of the actual
weight \(k\) modulo \(268\).  The continuous functions
\(x^{k-1}\eta(t)\) and \(x^{n-1}\eta(t)\) are congruent
modulo \(269\) on \emph{all} of the measure's domain.  On units
this is Fermat's theorem; on \(p\mathbb Z_p\), both exponents
are positive, so both powers reduce to zero.  Applying the
integral measure and taking constant coefficients gives
\[
                 c_k:=\frac{L(1-k,\eta)}2
                         \equiv c_n:=\frac{L(1-n,\eta)}2
                                        \pmod{269}.
 \tag{G116.18}
\]
The equal smoothing factors are units by (G116.17).
This remains valid when \(p\mid k\): no division by the actual
weight is made modulo \(p\).  The excluded residue \(n=1\)
is not treated by an invalid comparison at \(x=0\).

Here is the complete finite certificate for the \(133\) possible
reference weights.  In \(\mathbb F_{269}\), calculate
\[
\begin{split}
 B_0&=1,\qquad
 B_m=-\frac1{m+1}\sum_{i<m}\binom{m+1}{i}B_i
                                      \quad(1\le m\le267),\\
 B_n(X)&=\sum_{i=0}^n\binom ni B_iX^{n-i},\\
 c_n&=-\frac{3^{n-1}}{2n}
                   \bigl(B_n(1/3)-B_n(2/3)\bigr).
\end{split}
 \tag{G116.19}
\]
Every displayed denominator is prime to \(269\).  In increasing
order \(n=3,5,\ldots,267\), the resulting least nonnegative
residues of \(c_n\) are
\[
\begin{gathered}
239,90,177,20,12,152,229,165,214,184,35,204\\
133,199,187,66,72,171,220,198,265,131,217,52\\
67,99,183,252,184,214,14,91,110,106,103,219\\
120,104,228,214,231,141,219,56,31,214,97,142\\
168,178,103,164,25,248,173,46,122,90,138,4\\
104,261,79,126,180,93,7,79,45,78,178,3\\
66,251,151,162,245,253,27,108,253,262,86,8\\
209,203,160,152,122,239,76,129,143,174,27,215\\
120,109,262,128,235,141,216,138,92,203,27,254\\
54,63,66,80,203,222,133,90,257,257,234,85\\
118,242,150,17,191,8,12,74,142,72,89,42\\
71
\end{gathered}
 \tag{G116.20}
\]
All \(133\) entries are nonzero.  The retained script implements
exactly (G116.19), checks the length and each nonzero entry,
and records the generalized Bernoulli, \(L\)-value, constant
and smoothing factor for every reference weight.
A separate rational Akiyama--Tanigawa computation, with
reduction only after the exact rational generalized Bernoulli
calculation, agrees with every row.  No floating-point estimate,
GRH input, table lookup or incomplete class-group calculation
is used.  For example the last row gives
\(B_{267,\eta}=15\), \(L(1-267,\eta)=142\), and \(c_{267}=71\)
modulo \(269\).  Equations (G116.18)--(G116.20) therefore make the
actual constant \(c_k\) a unit for every retained weight.

\proofstep{6. The constant-expansion contradiction.}
Pull both forms to full level \(6\), and enlarge the finite
value field once to include the required roots of unity; write
\(O'\) for its integers.  The original deformation coefficient
ring, including its retained \(e_O=1\), is unchanged.
This is prime-to-\(p\) level.  Katz's integral generalized
modular-form ring \(V(O',\Gamma(6))\) admits arbitrary
coefficient base change and an injective expansion map at a
chosen cusp, by \S1.2 of the same paper.
The integral measure places \(E_k\) in this ring.
The form \(g\) also belongs to it: its integral expansion
and the divided-congruence inclusion in \S1.2 give this
directly.  Equivalently clear a fixed coefficient denominator
in the classical form and use the injective expansion map
modulo each power of the coefficient uniformizer to divide
back.  The arbitrary-ground-ring base-change statement makes
this argument valid for a ramified finite coefficient field.

Both forms have the same exact weight character \(a\mapsto a^k\).
After reduction modulo \(\mathfrak m_E\), (G116.15)--(G116.20)
say that \(E_k-g\) has the nonzero constant expansion
\(\overline c_k=\overline c_n\).  Injectivity identifies it
with that constant generalized form.  Applying
\(a=2\in\mathbb Z_p^\times\) gives simultaneously
\[
                        2^k\,\overline c_k=\overline c_k.
 \tag{G116.21}
\]
This is impossible: \(\overline c_k\ne0\) by (G116.20),
and \(2^k\ne1\) by (G116.17).

The same last step is precisely the \(m=1\) instance of
Katz, \emph{p-adic properties of modular schemes and modular
forms}, Corollary 4.4.1: a nonzero constant expansion of
weight \(k\) on an ordinary component forces \(p-1\mid k\).
Its use does not require \(p>k+1\), so there is no restriction
to a bounded range of weights.  The contradiction proves the
theorem.
\end{proof}

\begin{corollary}[AC-G116-seed: the original-endpoint seed instance is empty]
\label{thm:ACG116seed}
Under the original \(N_\psi=4\), \(\psi=\chi_{-4}\) datum and
all retained hypotheses of AC-G116-exclusion, an actual old Type II
point cannot supply a starting modular endpoint for the inherited
AC-G78 application, the automorphic instance in AC-G103-certificate,
or the native rank-two construction AC-G106-native.
This applies in particular to both weight residue classes \(133\) and \(267\)
of AC-G114--AC-G115.  The pure arithmetic strict-vanishing conclusions
of those two results remain independent of any such endpoint.
\end{corollary}
\begin{proof}
Each of the three named original-level applications explicitly
retains an actual old Type II modular endpoint under the old source
modularity input; none proves its existence from a residual tuple.
Apply the representation-level form of AC-G116-exclusion directly.
The native \texttt{typeII\_facts} in
\texttt{CoboundaryLemmaNoColumn.lean} constructs its odd endpoint
\(V\) with determinant \(\eta\).  Its displayed \texttt{Wfun}
formula and the original residual condition give (G116.3).
The original away-inertia and coefficient-prime flag conclusions,
together with the retained selected potentially crystalline local
condition, give the same local hypotheses used in Steps 1--3.
The native \texttt{admissible\_of\_typeII} argument, or its dense
variant \texttt{adm2B\_rhoK\_of\_typeII\_of\_modular\_extensions}
in \texttt{DenseAdmissibleOfTypeII.lean}, supplies the actual
newform under the retained old modularity input.  Thus this old
endpoint itself has exactly the properties already contradicted.
This argument needs no reverse old-to-cubic map and no claim that
an arbitrary old point belongs to the selected cubic prime.
Thus this particular supply of the seed is unavailable.
The algebraic implications proved conditionally from an endpoint
or from native support remain valid conditional statements.
The finite Kummer/unit proofs and original-uniformizer argument of
AC-G114--AC-G115 have no endpoint premise, so their nonvacuous
strict arithmetic conclusions are unaffected.

The generic AC-G78 construction permits an ordinary input at an
auxiliary tame level.  The exclusion does not apply merely because
that input has the same residual trace: its proof uses the exact
original scalar inertia at \(2\), quadratic inertia at \(3\), and absence
of other away ramification to obtain the characteristic-zero
primitive conductor \(3\) in (G116.6).  An auxiliary-level input must
independently supply its required local types and a legitimate
comparison with the target.  No existence of such an input, host,
or comparison is asserted by this corollary.
\end{proof}

In particular this section proves no original-H04.4 conclusion,
mixed Selmer vanishing, regulator nonvanishing, quantitative
counting bound, or unconditional instantiation of the endpoint-based
propagation pipeline.  It does not exclude a seed obtained by a
different construction with independently verified prerequisites.

\paragraph{Exact source scope.}
The native containing files \texttt{Newform.lean},
\texttt{ConjugateSelfDual/Statement.lean}, and
\texttt{AdmissibleOfTypeII.lean} were read in full, together
with the applied original-point arguments in AC-G68 and the
good-twist proof.  The exact \texttt{cycZ\_of\_isFrob} clause
was also checked.  No Lean build is claimed.

For local compatibility, the fresh primary checks are
Carayol, \emph{Sur les representations l-adiques associees
aux formes modulaires de Hilbert}, \S0.1--0.7, especially
Theorem A, printed pp.~409--411, at
\url{https://numdam.org/item/10.24033/asens.1512.pdf};
and Saito's own precise restatement of the classical theorem,
\emph{Galois representations in arithmetic geometry},
\S2a and \S2c, Theorem 3, printed pp.~6,8--9, at
\url{https://www.ms.u-tokyo.ac.jp/~t-saito/pp/suams.pdf}.
The latter explicitly uses geometric Frobenius in the local
Euler factor.  The original Saito theorem is
\emph{Inventiones} 129 (1997), pp.~607--620.
These deep compatibility and classical newvector theorems
are cited inputs; no new proof of them is claimed.
For the primitive level, Casselman's \emph{On some results of
Atkin and Lehner}, \emph{Mathematische Annalen} 201 (1973),
Theorem 1, printed p.~302, gives the local newvector conductor;
Theorem 4 and its preceding global application, printed
pp.~312--313, give the primitive global level as the product of
these local conductors.  Those statements and that global
application were checked directly on the scanned published pages
at \url{https://lesesvre.perso.math.cnrs.fr/newforms-references/casselman.pdf}.
The entire local proof of Theorem 1 is not claimed as a new source read.

For the Eisenstein step, Katz's primary paper is
\url{https://web.math.princeton.edu/~nmk/old/eismeasure.pdf}.
The exact definitions and base-change/expansion statements of
\S1.1--1.2, the complete printed Key Lemmas 1.2.1--1.2.3,
\S2.1.1--2.1.2, \S2.5 with the complete printed Lemma 2.5.3,
and \S3.1--3.3 including the complete printed proof of
Theorem 3.3.3 and Proposition 3.3.9 were read.
The moment exponent in (3.3.2) was also checked on the
rendered printed page 266.  The independent constant-expansion
source is
\url{https://web.math.princeton.edu/~nmk/old/padicpropMFMS.pdf},
\S4.2 and the complete printed proofs of Theorems 4.3/4.3bis
and Corollaries 4.4.1--4.4.2, printed pp.~149--152.
The foundational Igusa theorem and modular-curve integral
theory remain imported at that source scope.

\paragraph{Version and computational review scope.}
The earlier restricted diagnostic, SHA256
\texttt{6a848b771b0b64f4a815494e435725ade53f150b3272076032047c7d2feb93c6},
is retained unchanged with its author and independent review events.
The present theorem changes the weight scope and supplies the entire
finite certificate needed for that change; it is not a silent widening
of the earlier reviewed statement.  The exact original \texttt{GD}
containing file \texttt{Counts.lean}, the full AC-G68-point proof,
and the named starting-endpoint clauses were freshly reread for this
scope comparison.  The complete containing files
\texttt{CoboundaryLemmaNoColumn.lean} and
\texttt{DenseAdmissibleOfTypeII.lean} were also freshly read for
the direct old-endpoint consequence.  No Lean build or new proof of the imported old
modularity theorem is claimed.
